\section{Comparison to Large-Scale Agent Evaluations}

The empirical role of this section is deliberately narrow. We compare the qualitative predictions of the
preceding topology results to architecture-level trends reported in large-scale external evaluations of
agent systems \cite{kim2025scaling}. These reported metrics are not literal plug-in values of the
simplified theorem parameters. They are used here only as external consistency checks.

Kim et al.~\cite{kim2025scaling} report several architecture-level patterns that align with the topology
results in Section~5. Independent coordination exhibits much higher error amplification than centralized
coordination ($A_e = 17.2$ versus $4.4$), consistent with the claim that a validation bottleneck can
suppress unchecked worker errors, though these aggregates should not be read as literal estimates of the
theorem parameters. PlanCraft degrades under every multi-agent topology, from $-39.1\%$ to $-70.1\%$
relative to SAS, consistent with the sequential-penalty result that manufactured handoffs are pure
overhead when they add no new observations. Finance-Agent shows the opposite regime: centralized MAS
improves success from $0.349$ to $0.631$ ($+80.8\%$), which is consistent with the positive case for
decomposable tasks, but still does not directly estimate the speedup variable $S$. In the tool-heavy
Workbench domain, $T=16$ tools with typical $n=3$ agent configurations implies substantial remote-tool
use, and the reported negative efficiency--tool interaction
($\hat{\beta}_{E_c \times T} = -0.267$, $p < 0.001$) is consistent with the coordination tax from
the tool-density bound (Definition~\ref{thm:tool-density}).

Taken together, these results support the topology-level ordering rather than a numerical fit: reviewer
or hub structures help when multiple workers can independently introduce costly errors, strictly
sequential tasks should remain consolidated unless a handoff introduces a real new capability, parallel
specialist swarms help when subtasks are close to conditionally independent, and tool fragmentation must
be treated as an explicit routing problem rather than a free decomposition.

\begin{table}[h]
\centering
\small
\begin{tabular}{@{}>{\raggedright\arraybackslash}p{0.18\textwidth}>{\raggedright\arraybackslash}p{0.20\textwidth}>{\raggedright\arraybackslash}p{0.25\textwidth}>{\raggedright\arraybackslash}p{0.25\textwidth}@{}}
\toprule
\textbf{Topology} & \textbf{Epistemic Property} & \textbf{Prediction} & \textbf{Observed} \\
\midrule
Independent ($\otimes$) & No inter-agent $K_i$ & Highest error amplification & Architecture-level $A_e = 17.2$ \\
Centralized ($\circ$+hub) & Hub pools worker outputs & Validation bottleneck lowers amplification & Architecture-level $A_e = 4.4$ \\
Multi-agent + sequential & Requires manufactured $K_{\mathrm{receiver}}$ & Strictly sequential tasks degrade under handoffs & PlanCraft: $-39.1\%$ to $-70.1\%$ vs.\ SAS \\
Multi-agent + parallel & Approximate epistemic independence & Large gains on decomposable tasks & Finance-Agent Centralized: $+80.8\%$ vs.\ SAS \\
Multi-agent + high $|T|$ & Knowledge fragmentation & Coordination tax grows with $t$ and $n$ & Workbench $T=16$; $\hat{\beta}_{E_c \times T} = -0.267$ \\
\bottomrule
\end{tabular}
\caption{Epistemic predictions versus empirical observations from Kim et al.~\cite{kim2025scaling}. The
table compares topology-level qualitative predictions to architecture-level metrics. Reported
percentages and error-amplification factors are empirical aggregates, not literal plug-in values of the
simplified theorem parameters.}
\label{tab:epistemic-predictions}
\end{table}

This comparison is limited to a single external study~\cite{kim2025scaling}.
Broader empirical validation across multiple benchmarks and framework
configurations remains future work.
